\documentclass[12pt]{ximera}
\typeout{Start loading xmPreamble.tex}%

% Add here extra macro's that are loaded automatically by all documents of claas 'ximera' or 'xourse' in this repo

%%
%%  Example:
%%
% \newcommand{\R}{\mathbb{R}}
\usepackage{tikz}
\usetikzlibrary{positioning, arrows.meta}
\usepackage{multicol}
\usepackage{enumitem}
\usepackage{caption}
\usepackage{amsmath}
\usepackage{amsthm}
\usepackage{fontenc}

%% ---------------------------------------------------------------
%% Shared worksheet macros.
%%
%% These came from the Summer 2026 worksheet set, where each worksheet carried its
%% own copy. They have to live here instead: a xourse inlines every activity into a
%% single document, so duplicate \newcommand definitions across activities collide,
%% and the xourse gobbles \input so a per-topic macro file never loads.
%%
%% They use \providecommand, NOT \newcommand. ximera.cls loads this file automatically,
%% and every activity also carries an explicit \typeout{Start loading xmPreamble.tex}%

% Add here extra macro's that are loaded automatically by all documents of claas 'ximera' or 'xourse' in this repo

%%
%%  Example:
%%
% \newcommand{\R}{\mathbb{R}}
\usepackage{tikz}
\usetikzlibrary{positioning, arrows.meta}
\usepackage{multicol}
\usepackage{enumitem}
\usepackage{caption}
\usepackage{amsmath}
\usepackage{amsthm}
\usepackage{fontenc}

%% ---------------------------------------------------------------
%% Shared worksheet macros.
%%
%% These came from the Summer 2026 worksheet set, where each worksheet carried its
%% own copy. They have to live here instead: a xourse inlines every activity into a
%% single document, so duplicate \newcommand definitions across activities collide,
%% and the xourse gobbles \input so a per-topic macro file never loads.
%%
%% They use \providecommand, NOT \newcommand. ximera.cls loads this file automatically,
%% and every activity also carries an explicit \typeout{Start loading xmPreamble.tex}%

% Add here extra macro's that are loaded automatically by all documents of claas 'ximera' or 'xourse' in this repo

%%
%%  Example:
%%
% \newcommand{\R}{\mathbb{R}}
\usepackage{tikz}
\usetikzlibrary{positioning, arrows.meta}
\usepackage{multicol}
\usepackage{enumitem}
\usepackage{caption}
\usepackage{amsmath}
\usepackage{amsthm}
\usepackage{fontenc}

%% ---------------------------------------------------------------
%% Shared worksheet macros.
%%
%% These came from the Summer 2026 worksheet set, where each worksheet carried its
%% own copy. They have to live here instead: a xourse inlines every activity into a
%% single document, so duplicate \newcommand definitions across activities collide,
%% and the xourse gobbles \input so a per-topic macro file never loads.
%%
%% They use \providecommand, NOT \newcommand. ximera.cls loads this file automatically,
%% and every activity also carries an explicit \input{../xmPreamble}, so this file is
%% read twice. That was harmless while it held only \usepackage lines (LaTeX skips a
%% package it has already loaded) but a second \newcommand is a hard error.
%%
%%   \blank[width]        fill-in-the-blank rule (default 2.5cm)
%%   \showwork            "show and explain your work" reminder
%%   \showworkline        ... plus which schedule line was used
%%   \showworkyear        ... plus which year's schedule was used
%%   \schedhead{...}      centred bold caption above a tiered-rate table
%%   \samesquares[scale]{left}{right}   two equal-area squares, labelled
%%   \samecubes[scale]{left}{right}     two equal-volume cubes, labelled
%% ---------------------------------------------------------------

\providecommand{\blank}[1][2.5cm]{\underline{\hspace{#1}}}
\providecommand{\showwork}{\textbf{REMEMBER TO SHOW AND EXPLAIN YOUR WORK.}}
\providecommand{\showworkline}{\textbf{REMEMBER TO SHOW AND EXPLAIN YOUR WORK.
  Explanation should include how and why you know which line of the
  schedule was used to calculate this amount.}}
\providecommand{\showworkyear}{\begin{sloppypar}\textbf{REMEMBER TO SHOW AND
  EXPLAIN YOUR WORK. Explanation should include how and why you know which
  line of the year's tax schedule was used to calculate this tax
  amount.}\end{sloppypar}}
\providecommand{\schedhead}[1]{\begin{center}\textbf{#1}\end{center}}

\providecommand{\samesquares}[3][1]{%
\begin{center}
{\footnotesize These squares are the \textbf{same size}.}

\vspace{1.5mm}
\begin{tikzpicture}[scale=#1,every node/.style={scale=#1}]
  \draw (0,0) rectangle (2.4,2.4);
  \node[anchor=north west] at (0.15,2.25) {Area:};
  \node[anchor=east]  at (-0.15,1.2) {#2};
  \node[anchor=north] at (1.2,-0.15) {#2};
  \begin{scope}[xshift=4.6cm]
    \draw (0,0) rectangle (2.4,2.4);
    \node[anchor=north west] at (0.15,2.25) {Area:};
    \node[anchor=east]  at (-0.15,1.2) {#3};
    \node[anchor=north] at (1.2,-0.15) {#3};
  \end{scope}
\end{tikzpicture}
\end{center}}

\providecommand{\samecubes}[3][1]{%
\begin{center}
{\footnotesize These cubes are the \textbf{same size}.}

\vspace{1.5mm}
\begin{tikzpicture}[scale=#1,every node/.style={scale=#1}]
  \foreach \dx in {0,5.2} {
    \begin{scope}[xshift=\dx cm]
      \draw (0,0) rectangle (2.0,2.0);
      \draw (0,2.0) -- (0.65,2.6) -- (2.65,2.6) -- (2.0,2.0);
      \fill[black!12] (2.0,0) -- (2.65,0.6) -- (2.65,2.6) -- (2.0,2.0) -- cycle;
      \draw (2.0,0) -- (2.65,0.6) -- (2.65,2.6) -- (2.0,2.0);
      \node[anchor=north west] at (0.1,1.9) {Volume:};
    \end{scope}
  }
  \node[anchor=east]  at (-0.15,1.0) {#2};
  \node[anchor=north] at (1.0,-0.15) {#2};
  \node[anchor=west]  at (2.25,0.4) {#2};
  \node[anchor=east]  at (5.05,1.0) {#3};
  \node[anchor=north] at (6.2,-0.15) {#3};
  \node[anchor=west]  at (7.45,0.4) {#3};
\end{tikzpicture}
\end{center}}

%% NOTE: do NOT add \usepackage to individual activity .tex files.
%% When a xourse inlines an activity it gobbles \documentclass and \input, but a bare
%% \usepackage still executes -- after \begin{document} -- and errors out.
%% All shared packages and macros belong here.
%%
%% ximera.cls already loads: amsmath, amssymb, amsthm, cancel, enumitem, graphicx,
%% hyperref, listings, pgfplots, tikz, url, xcolor. No need to re-load those.
, so this file is
%% read twice. That was harmless while it held only \usepackage lines (LaTeX skips a
%% package it has already loaded) but a second \newcommand is a hard error.
%%
%%   \blank[width]        fill-in-the-blank rule (default 2.5cm)
%%   \showwork            "show and explain your work" reminder
%%   \showworkline        ... plus which schedule line was used
%%   \showworkyear        ... plus which year's schedule was used
%%   \schedhead{...}      centred bold caption above a tiered-rate table
%%   \samesquares[scale]{left}{right}   two equal-area squares, labelled
%%   \samecubes[scale]{left}{right}     two equal-volume cubes, labelled
%% ---------------------------------------------------------------

\providecommand{\blank}[1][2.5cm]{\underline{\hspace{#1}}}
\providecommand{\showwork}{\textbf{REMEMBER TO SHOW AND EXPLAIN YOUR WORK.}}
\providecommand{\showworkline}{\textbf{REMEMBER TO SHOW AND EXPLAIN YOUR WORK.
  Explanation should include how and why you know which line of the
  schedule was used to calculate this amount.}}
\providecommand{\showworkyear}{\begin{sloppypar}\textbf{REMEMBER TO SHOW AND
  EXPLAIN YOUR WORK. Explanation should include how and why you know which
  line of the year's tax schedule was used to calculate this tax
  amount.}\end{sloppypar}}
\providecommand{\schedhead}[1]{\begin{center}\textbf{#1}\end{center}}

\providecommand{\samesquares}[3][1]{%
\begin{center}
{\footnotesize These squares are the \textbf{same size}.}

\vspace{1.5mm}
\begin{tikzpicture}[scale=#1,every node/.style={scale=#1}]
  \draw (0,0) rectangle (2.4,2.4);
  \node[anchor=north west] at (0.15,2.25) {Area:};
  \node[anchor=east]  at (-0.15,1.2) {#2};
  \node[anchor=north] at (1.2,-0.15) {#2};
  \begin{scope}[xshift=4.6cm]
    \draw (0,0) rectangle (2.4,2.4);
    \node[anchor=north west] at (0.15,2.25) {Area:};
    \node[anchor=east]  at (-0.15,1.2) {#3};
    \node[anchor=north] at (1.2,-0.15) {#3};
  \end{scope}
\end{tikzpicture}
\end{center}}

\providecommand{\samecubes}[3][1]{%
\begin{center}
{\footnotesize These cubes are the \textbf{same size}.}

\vspace{1.5mm}
\begin{tikzpicture}[scale=#1,every node/.style={scale=#1}]
  \foreach \dx in {0,5.2} {
    \begin{scope}[xshift=\dx cm]
      \draw (0,0) rectangle (2.0,2.0);
      \draw (0,2.0) -- (0.65,2.6) -- (2.65,2.6) -- (2.0,2.0);
      \fill[black!12] (2.0,0) -- (2.65,0.6) -- (2.65,2.6) -- (2.0,2.0) -- cycle;
      \draw (2.0,0) -- (2.65,0.6) -- (2.65,2.6) -- (2.0,2.0);
      \node[anchor=north west] at (0.1,1.9) {Volume:};
    \end{scope}
  }
  \node[anchor=east]  at (-0.15,1.0) {#2};
  \node[anchor=north] at (1.0,-0.15) {#2};
  \node[anchor=west]  at (2.25,0.4) {#2};
  \node[anchor=east]  at (5.05,1.0) {#3};
  \node[anchor=north] at (6.2,-0.15) {#3};
  \node[anchor=west]  at (7.45,0.4) {#3};
\end{tikzpicture}
\end{center}}

%% NOTE: do NOT add \usepackage to individual activity .tex files.
%% When a xourse inlines an activity it gobbles \documentclass and \input, but a bare
%% \usepackage still executes -- after \begin{document} -- and errors out.
%% All shared packages and macros belong here.
%%
%% ximera.cls already loads: amsmath, amssymb, amsthm, cancel, enumitem, graphicx,
%% hyperref, listings, pgfplots, tikz, url, xcolor. No need to re-load those.
, so this file is
%% read twice. That was harmless while it held only \usepackage lines (LaTeX skips a
%% package it has already loaded) but a second \newcommand is a hard error.
%%
%%   \blank[width]        fill-in-the-blank rule (default 2.5cm)
%%   \showwork            "show and explain your work" reminder
%%   \showworkline        ... plus which schedule line was used
%%   \showworkyear        ... plus which year's schedule was used
%%   \schedhead{...}      centred bold caption above a tiered-rate table
%%   \samesquares[scale]{left}{right}   two equal-area squares, labelled
%%   \samecubes[scale]{left}{right}     two equal-volume cubes, labelled
%% ---------------------------------------------------------------

\providecommand{\blank}[1][2.5cm]{\underline{\hspace{#1}}}
\providecommand{\showwork}{\textbf{REMEMBER TO SHOW AND EXPLAIN YOUR WORK.}}
\providecommand{\showworkline}{\textbf{REMEMBER TO SHOW AND EXPLAIN YOUR WORK.
  Explanation should include how and why you know which line of the
  schedule was used to calculate this amount.}}
\providecommand{\showworkyear}{\begin{sloppypar}\textbf{REMEMBER TO SHOW AND
  EXPLAIN YOUR WORK. Explanation should include how and why you know which
  line of the year's tax schedule was used to calculate this tax
  amount.}\end{sloppypar}}
\providecommand{\schedhead}[1]{\begin{center}\textbf{#1}\end{center}}

\providecommand{\samesquares}[3][1]{%
\begin{center}
{\footnotesize These squares are the \textbf{same size}.}

\vspace{1.5mm}
\begin{tikzpicture}[scale=#1,every node/.style={scale=#1}]
  \draw (0,0) rectangle (2.4,2.4);
  \node[anchor=north west] at (0.15,2.25) {Area:};
  \node[anchor=east]  at (-0.15,1.2) {#2};
  \node[anchor=north] at (1.2,-0.15) {#2};
  \begin{scope}[xshift=4.6cm]
    \draw (0,0) rectangle (2.4,2.4);
    \node[anchor=north west] at (0.15,2.25) {Area:};
    \node[anchor=east]  at (-0.15,1.2) {#3};
    \node[anchor=north] at (1.2,-0.15) {#3};
  \end{scope}
\end{tikzpicture}
\end{center}}

\providecommand{\samecubes}[3][1]{%
\begin{center}
{\footnotesize These cubes are the \textbf{same size}.}

\vspace{1.5mm}
\begin{tikzpicture}[scale=#1,every node/.style={scale=#1}]
  \foreach \dx in {0,5.2} {
    \begin{scope}[xshift=\dx cm]
      \draw (0,0) rectangle (2.0,2.0);
      \draw (0,2.0) -- (0.65,2.6) -- (2.65,2.6) -- (2.0,2.0);
      \fill[black!12] (2.0,0) -- (2.65,0.6) -- (2.65,2.6) -- (2.0,2.0) -- cycle;
      \draw (2.0,0) -- (2.65,0.6) -- (2.65,2.6) -- (2.0,2.0);
      \node[anchor=north west] at (0.1,1.9) {Volume:};
    \end{scope}
  }
  \node[anchor=east]  at (-0.15,1.0) {#2};
  \node[anchor=north] at (1.0,-0.15) {#2};
  \node[anchor=west]  at (2.25,0.4) {#2};
  \node[anchor=east]  at (5.05,1.0) {#3};
  \node[anchor=north] at (6.2,-0.15) {#3};
  \node[anchor=west]  at (7.45,0.4) {#3};
\end{tikzpicture}
\end{center}}

%% NOTE: do NOT add \usepackage to individual activity .tex files.
%% When a xourse inlines an activity it gobbles \documentclass and \input, but a bare
%% \usepackage still executes -- after \begin{document} -- and errors out.
%% All shared packages and macros belong here.
%%
%% ximera.cls already loads: amsmath, amssymb, amsthm, cancel, enumitem, graphicx,
%% hyperref, listings, pgfplots, tikz, url, xcolor. No need to re-load those.

\title{Probability}
\begin{document}
\begin{abstract}
Probability as a count of favourable outcomes over total outcomes, complementary
events, and the conditional probabilities that two-way tables make visible --- including
why ``the chance a ticketed driver drives red'' is a different question from ``the chance
a red-car driver gets ticketed.'' Then counting techniques, permutations versus
combinations, experimental versus theoretical probability, and dependent versus
independent events.
\end{abstract}
\maketitle

\section*{Introduction to Probability}

Probability is the likelihood of some instance or outcome occurring. To know how likely
something is to occur, we need to know how many total options there are and then how
many of those options fit the ``event'' --- what we call the instance we are looking to
occur.

Consider a dice-rolling scenario. The likelihood of rolling a ``1'' on a 6-sided die is
$1/6$, as there are 6 total sides and only one of them is a ``1.'' So there is 1 option
out of the 6 total options that fits the event we are looking for.

This assumes all 6 options are equally likely. If the sides 3 and 4 were more likely to
land face up for some reason --- a weighted or unfair die --- we would have to weight
those probabilities higher, much as we did in the Means \& Medians activity with the
averages of Wing Aero employees. You may assume all options are equally possible in the
problems here.

While we are often calculating a fraction or ratio for the probability, we prefer to
write probabilities as percentages. So the likelihood of rolling a 1 on a 6-sided die is
$1/6 = 0.1667$, which is a $16.67\%$ likelihood of occurring.

What does that $16.67\%$ mean? It means that if we rolled a die 100 times, we would
expect a ``1'' about 16.67 times. This does not mean it \emph{will} roll a ``1'' that
many times --- in fact it cannot, since you can't have $0.67$ of a roll. But it gives us
an idea of how likely something is, and expressing probabilities out of 100 lets us
compare likelihoods against each other.

\begin{problem}
Find the probability that a randomly drawn card from a deck of cards is a King. Give
your answer as a percentage rounded to two decimal places.

$\answer[tolerance=0.1]{7.69}\%$

\begin{solution}
There are 4 Kings in a standard deck of 52 cards, so
\[
\frac{4}{52} = \frac{1}{13} = 0.0769\ldots \approx 7.69\%.
\]
\end{solution}
\end{problem}

\begin{problem}
You have a bag containing 24 apples: eight are green and sixteen are red. If you pick an
apple randomly, what is the probability that it will be red?

$\answer[tolerance=0.1]{66.67}\%$

\begin{solution}
16 of the 24 apples are red, so
\[
\frac{16}{24} = \frac{2}{3} = 0.6667 \approx 66.67\%.
\]
\end{solution}
\end{problem}

\begin{problem}
A 6-sided die is rolled once. What is the probability that you see a number bigger than
4?

$\answer[tolerance=0.1]{33.33}\%$

\begin{solution}
The numbers bigger than 4 are 5 and 6 --- two outcomes out of six:
\[
\frac{2}{6} = \frac{1}{3} = 0.3333 \approx 33.33\%.
\]
Note that ``bigger than 4'' excludes 4 itself. If the question had said ``4 or bigger''
the answer would be $3/6 = 50\%$.
\end{solution}
\end{problem}

\section*{Complementary Events}

Complementary events are what we might think of as ``opposite'' events. There is a 1 in
6 chance of rolling a ``1,'' which means there is a 5 out of 6 chance, $5/6 = 0.8333$, of
NOT rolling a ``1.''

The mathematical relationship between complementary events is that their probabilities
always add to $100\%$. In our example, $1/6 + 5/6 = 16.67\% + 83.33\% = 100\%$.

We can use this to our advantage. It is sometimes easier to calculate the complement of
an event than the event itself --- if we know the probability of the complement, we
subtract it from $100\%$ to find the probability of the original event.

\begin{problem}
The Ohio State football team had a preseason probability of $52.4\%$ to go over 10.5 wins
in 2023. What was the preseason probability of going under 10.5 wins?

$\answer[tolerance=0.1]{47.6}\%$

\begin{solution}
Going over 10.5 wins and going under 10.5 wins are complementary --- with a half-win
threshold there is no possibility of landing exactly on it, so one or the other must
happen. Therefore
\[
100\% - 52.4\% = 47.6\%.
\]
\end{solution}
\end{problem}

\begin{problem}
When two cards are drawn from a deck randomly without replacement, the probability of
drawing a red card and an ace is $1/13$. What would be the complement to this event, and
what would be its likelihood of occurring? Give your answer as a percentage.

$\answer[tolerance=0.1]{92.31}\%$

\begin{freeResponse}
\end{freeResponse}

\begin{solution}
The complement is the event ``NOT drawing a red card and an ace'' --- that is, any
outcome other than getting both a red card and an ace. Its probability is
\[
100\% - \frac{1}{13} = 100\% - 7.69\% = 92.31\%.
\]

Be careful with the wording here: the complement of ``a red card and an ace'' is not
``a black card and a non-ace.'' The complement of an event is \emph{everything else},
which includes drawing a red card without an ace, an ace that isn't red, and neither.
\end{solution}
\end{problem}

\section*{Two-Way Tables and Conditional Probability}

The table shows the number of drivers who received and did not receive a speeding ticket
on a college campus over a three year period, and the colour of their car.

\begin{center}
\begin{tabular}{lccc}
\hline
 & Received ticket & Did not receive ticket & Total \\
\hline
Red car & 18 & 162 & 180 \\
Not a red car & 54 & 564 & 618 \\
\hline
Total & 72 & 726 & 798 \\
\hline
\end{tabular}
\end{center}

\begin{problem}
Determine the probability that a driver has a red car AND gets a speeding ticket. Write
a sentence about what this percentage means.

$\answer[tolerance=0.1]{2.26}\%$

\begin{freeResponse}
\end{freeResponse}

\begin{solution}
18 drivers are in both categories, out of 798 drivers total:
\[
\frac{18}{798} = 0.02256 \approx 2.26\%.
\]
Meaning: of all the drivers on campus, about $2.26\%$ both drive a red car and got a
ticket. The denominator is everybody.
\end{solution}
\end{problem}

\begin{problem}
Determine the probability that a driver has a red car OR gets a speeding ticket. Write a
sentence about what this percentage means.

$\answer[tolerance=0.1]{29.32}\%$

\begin{freeResponse}
\end{freeResponse}

\begin{solution}
Add the red-car drivers and the ticketed drivers, then subtract the overlap so it is not
counted twice:
\[
\frac{180 + 72 - 18}{798} = \frac{234}{798} = 0.29323 \approx 29.32\%.
\]
Meaning: about $29.32\%$ of drivers fall into at least one of the two categories.

The subtraction matters. Without it you would get $\frac{252}{798} = 31.58\%$, which
double-counts the 18 drivers who are in both groups.
\end{solution}
\end{problem}

\begin{problem}
Determine the probability that a driver \emph{of a red car} gets a speeding ticket.
Write a sentence about what this percentage means.

$\answer[tolerance=0.1]{10}\%$

\begin{freeResponse}
\end{freeResponse}

\begin{solution}
Now we are only looking at red-car drivers, so the denominator is 180, not 798:
\[
\frac{18}{180} = 0.10 = 10\%.
\]
Meaning: among drivers of red cars, $10\%$ got a ticket.
\end{solution}
\end{problem}

\begin{problem}
Determine the probability that a driver \emph{who gets a speeding ticket} has a red car.
Write a sentence about what this percentage means.

$\answer[tolerance=0.1]{25}\%$

\begin{freeResponse}
\end{freeResponse}

\begin{solution}
Now we are only looking at ticketed drivers, so the denominator is 72:
\[
\frac{18}{72} = 0.25 = 25\%.
\]
Meaning: among drivers who got a ticket, $25\%$ drove a red car.

Compare this with the previous problem. Both use the same 18 drivers in the numerator,
but they answer completely different questions because the denominators differ --- $10\%$
of red-car drivers get ticketed, while $25\%$ of ticketed drivers drive red. Reversing
which group you condition on can change the answer dramatically, and confusing the two is
one of the most common errors in reading statistics.
\end{solution}
\end{problem}

\begin{problem}
How likely are you to get a speeding ticket while driving a red car, as opposed to
another colour of car?

\begin{multipleChoice}
\choice{Significantly more likely}
\choice[correct]{Slightly more likely}
\choice{Equally likely}
\choice{Slightly less likely}
\choice{Significantly less likely}
\end{multipleChoice}

Support your choice by referencing two of the probabilities you calculated, and explain
why you chose those two.

\begin{freeResponse}
\end{freeResponse}

\begin{solution}
The two probabilities to compare are the chance of a ticket \emph{given} a red car and
the chance of a ticket \emph{given} a non-red car:
\[
P(\text{ticket} \mid \text{red}) = \frac{18}{180} = 10\%,
\qquad
P(\text{ticket} \mid \text{not red}) = \frac{54}{618} = 8.74\%.
\]
$10\%$ against $8.74\%$ is a real difference but a small one, so \emph{slightly more
likely} is the right answer.

Why these two? Because the question asks about your risk of a ticket \emph{conditional on
what colour car you drive}, so car colour has to be the thing held fixed in the
denominator. The $25\%$ figure --- the share of ticketed drivers with red cars --- looks
alarming but cannot answer this question: red cars could be $25\%$ of tickets simply by
being common on campus. And the $2.26\%$ figure has everybody in the denominator, so it
does not isolate colour at all.
\end{solution}
\end{problem}

\section*{Low Birthweight}

Suppose that at a local hospital there were 75 births last month, and 18 of those were to
women who smoked. Suppose further that 3 of the 18 babies born to mothers who smoked were
low birthweight, and 6 of the 57 babies born to mothers who did not smoke were low
birthweight.

\begin{problem}
Determine the probability of a low birthweight from a smoking mother.

$\answer[tolerance=0.1]{16.67}\%$

\begin{solution}
\[
\frac{3}{18} = 0.1667 \approx 16.67\%.
\]
\end{solution}
\end{problem}

\begin{problem}
Determine the probability that the mother of a low birthweight baby smokes.

$\answer[tolerance=0.1]{33.33}\%$

\begin{solution}
There are $3 + 6 = 9$ low birthweight babies in total, 3 of whom had mothers who smoked:
\[
\frac{3}{9} = 0.3333 \approx 33.33\%.
\]
\end{solution}
\end{problem}

\begin{problem}
Determine the probability of a low birthweight from a nonsmoking mother.

$\answer[tolerance=0.1]{10.53}\%$

\begin{solution}
\[
\frac{6}{57} = 0.10526 \approx 10.53\%.
\]
\end{solution}
\end{problem}

\begin{problem}
Determine the probability that the mother of a low birthweight baby does not smoke.

$\answer[tolerance=0.1]{66.67}\%$

\begin{solution}
\[
\frac{6}{9} = 0.6667 \approx 66.67\%.
\]
Note this and the earlier $33.33\%$ are complements --- every low birthweight baby has a
mother who either smoked or did not.
\end{solution}
\end{problem}

Use the probabilities above to mark each of the following True or False, then write a
sentence supporting your choice.

\begin{problem}
A baby born to a non-smoking mother is more likely to have a low birthweight than a baby
born to a smoking mother.
\begin{multipleChoice}
\choice{True}
\choice[correct]{False}
\end{multipleChoice}

\begin{freeResponse}
\end{freeResponse}

\begin{solution}
False. $P(\text{low BW} \mid \text{nonsmoker}) = 10.53\%$, which is lower than
$P(\text{low BW} \mid \text{smoker}) = 16.67\%$. The statement has it backwards.
\end{solution}
\end{problem}

\begin{problem}
A randomly selected low birthweight baby has a one-in-three chance of having a mother who
smokes.
\begin{multipleChoice}
\choice[correct]{True}
\choice{False}
\end{multipleChoice}

\begin{freeResponse}
\end{freeResponse}

\begin{solution}
True. $P(\text{smoker} \mid \text{low BW}) = 3/9 = 1/3 = 33.33\%$, which is exactly one
in three.
\end{solution}
\end{problem}

\begin{problem}
The majority of babies born to smoking mothers have a low birthweight.
\begin{multipleChoice}
\choice{True}
\choice[correct]{False}
\end{multipleChoice}

\begin{freeResponse}
\end{freeResponse}

\begin{solution}
False. Only $3$ of the $18$ babies born to smoking mothers were low birthweight ---
$16.67\%$, nowhere near a majority. The great majority, $83.33\%$, were not.
\end{solution}
\end{problem}

\begin{problem}
If a baby has a low birthweight, then the mother is more likely to be a smoker than a
non-smoker.
\begin{multipleChoice}
\choice{True}
\choice[correct]{False}
\end{multipleChoice}

\begin{freeResponse}
\end{freeResponse}

\begin{solution}
False. Among low birthweight babies, $33.33\%$ had mothers who smoked and $66.67\%$ had
mothers who did not. So the mother is twice as likely to be a non-smoker.

This is worth sitting with, because smoking does raise the risk. The reason non-smokers
still dominate among low birthweight babies is that there are far more of them --- 57
versus 18. A higher rate applied to a much smaller group can still produce fewer cases.
\end{solution}
\end{problem}

\begin{problem}
Most low birthweight babies are born to non-smoking mothers.
\begin{multipleChoice}
\choice[correct]{True}
\choice{False}
\end{multipleChoice}

\begin{freeResponse}
\end{freeResponse}

\begin{solution}
True --- $6$ of the $9$ low birthweight babies, or $66.67\%$, had non-smoking mothers,
for exactly the reason described above.
\end{solution}
\end{problem}

\begin{problem}
A baby born to a smoking mother is more likely to have a low birthweight than a baby born
to a non-smoking mother.
\begin{multipleChoice}
\choice[correct]{True}
\choice{False}
\end{multipleChoice}

\begin{freeResponse}
\end{freeResponse}

\begin{solution}
True. $16.67\%$ against $10.53\%$ --- smoking mothers' babies are more than one and a
half times as likely to be low birthweight.

Notice that this statement and the previous one are both true, even though they sound
like they contradict each other. Smoking raises an individual baby's risk, \emph{and}
most low birthweight babies still come from non-smoking mothers. Which is true depends
entirely on which group sits in the denominator.
\end{solution}
\end{problem}

\section*{Counting}

Amina is choosing a meal from a menu. She first decides between a soup or a salad. Then
she chooses a main dish: a sandwich, quiche, or fajita. Finally she selects either chips
or cookies as her side. Below is a decision tree listing all the possibilities.

\begin{center}
\includegraphics[width=0.85\textwidth]{probability/decision_tree.png}
\end{center}

\begin{problem}
How many possible meals does Amina have to choose from?

$\answer{12}$

\begin{solution}
Two starters, three mains, and two sides: $2 \times 3 \times 2 = 12$ meals, which is what
the tree shows.
\end{solution}
\end{problem}

\begin{problem}
Of these options, how many have fajita as part of the meal?

$\answer{4}$

\begin{solution}
Fix the main as fajita. The starter still has 2 options and the side still has 2, so
$2 \times 2 = 4$ meals include a fajita.
\end{solution}
\end{problem}

\begin{problem}
If Amina were to make her choices randomly, what would be the percentage likelihood of
having a fajita as part of her meal?

$\answer[tolerance=0.1]{33.33}\%$

\begin{solution}
\[
\frac{4}{12} = \frac{1}{3} \approx 33.33\%.
\]
Which makes sense: the fajita is one of three equally likely mains, and the other two
choices don't affect it.
\end{solution}
\end{problem}

Counting is important in determining probability --- you need to know how many total
options there are and how many are part of a certain outcome. Sometimes the things we
want to count are very large, so a technique helps.

In the decision tree there were two starter options, three main options, and two side
options. For each of the two starters there are three mains, making $2 \times 3$
combinations; and with each of those there are two sides, giving $2 \times 3 \times 2$.

In general, when we have a sequence of choices and we know how many options each has, we
find the total by \textbf{multiplying the options together}. If we build a meal from 3
fruits, 4 proteins, and 2 desserts, that is $3 \times 4 \times 2 = 24$ total combinations.

\begin{problem}
How many different ice cream sundaes can be made with the following choices: cup or cone,
two scoops of either strawberry, chocolate, or vanilla, and a topping of sprinkles,
M\&Ms, or fudge?

$\answer{54}$

\begin{freeResponse}
\end{freeResponse}

\begin{solution}
Treating each scoop as its own independent choice from the three flavours:
\[
\underbrace{2}_{\text{cup/cone}} \times \underbrace{3}_{\text{scoop 1}}
\times \underbrace{3}_{\text{scoop 2}} \times \underbrace{3}_{\text{topping}} = 54.
\]

This one depends on how you read the wording. If you took ``two scoops of either
strawberry, chocolate, or vanilla'' to mean both scoops must be the \emph{same} flavour,
you would get $2 \times 3 \times 3 = 18$. If you decided that a chocolate-then-vanilla
sundae is the same as vanilla-then-chocolate, you would get $2 \times 6 \times 3 = 36$.
All three readings are defensible; what matters is stating which one you used.
\end{solution}
\end{problem}

\begin{problem}
How many possible outfits can be made from 9 shirts, 3 pants, and 10 socks?

$\answer{270}$

\begin{solution}
\[
9 \times 3 \times 10 = 270.
\]
\end{solution}
\end{problem}

\begin{problem}
A ``build your own sandwich'' restaurant wants an advertisement saying ``Over \_\_\_\_\_
sandwich options!'' A customer chooses one from 3 breads, one from 4 proteins,
\emph{two} from 11 toppings, and one from 7 sauces. How many different sandwiches can be
made?

$\answer{4620}$

\begin{freeResponse}
\end{freeResponse}

\begin{solution}
The bread, protein, and sauce are straightforward: 3, 4, and 7 options. The toppings are
different, because you choose \emph{two} from 11 and the order you name them in does not
change the sandwich. That is a combination:
\[
\binom{11}{2} = \frac{11 \times 10}{2} = 55.
\]
So the total is
\[
3 \times 4 \times 55 \times 7 = 4{,}620.
\]

For the sign, you would round \emph{down} to a number you can honestly claim to be over
--- ``Over 4,500 sandwich options!'' is true, whereas ``Over 5,000'' would not be.

Note that this problem needs the combination idea from the next section. If you instead
treated the two toppings as ordered, you would get $11 \times 10 = 110$ and a total of
$9{,}240$ --- exactly double, because every topping pair would be counted twice.
\end{solution}
\end{problem}

\section*{Permutations \& Combinations}

Suppose you forget the passcode to a keypad with buttons 0 through 9, but you remember it
is 4 digits long and every digit was different. How many possible passcodes are there?

Or maybe you need to choose 4 of a group of 10 friends for an outing. How many different
options are there?

In each case we have a group of ``objects'' to choose from and a ``sample'' to choose from
that group. The difference is that in one case the \emph{order matters} --- 1567 is a
different passcode from 5671 --- and in the other it does not, since choosing Billy,
Jenny, Lily, and Fred is the same as choosing Fred, Lily, Billy, and Jenny.

The passcode is a \textbf{permutation}: a sample of 4 from 10 objects where order matters,
giving 5,040 possible passcodes. The friend group is a \textbf{combination}: a sample of 4
from 10 where order does not matter, giving 210 possible choices.

For each scenario below, decide whether a permutation or a combination is appropriate,
then calculate the answer.

\begin{problem}
A student has 5 different books and wants to arrange 3 of them in a row. How many
different ways can the books be arranged?
\begin{multipleChoice}
\choice[correct]{Permutation}
\choice{Combination}
\end{multipleChoice}

$\answer{60}$

\begin{solution}
Permutation --- ``arrange in a row'' means the order on the shelf matters.
$P(5,3) = 5 \times 4 \times 3 = 60$.
\end{solution}
\end{problem}

\begin{problem}
A group of 7 friends need to get in line for a ride. How many possible ways can they line
themselves up?
\begin{multipleChoice}
\choice[correct]{Permutation}
\choice{Combination}
\end{multipleChoice}

$\answer{5040}$

\begin{solution}
Permutation --- a line has an order. Here the whole group is being arranged, so
$P(7,7) = 7! = 5{,}040$.
\end{solution}
\end{problem}

\begin{problem}
A pizza shop offers 8 toppings. If you can choose any 3 toppings, how many different ways
can you order pizza?
\begin{multipleChoice}
\choice{Permutation}
\choice[correct]{Combination}
\end{multipleChoice}

$\answer{56}$

\begin{solution}
Combination --- the pizza is the same regardless of which order you named the toppings.
$\binom{8}{3} = 56$.
\end{solution}
\end{problem}

\begin{problem}
Out of a group of 6 friends, how many ways can you choose 2 to go to a concert with you?
\begin{multipleChoice}
\choice{Permutation}
\choice[correct]{Combination}
\end{multipleChoice}

$\answer{15}$

\begin{solution}
Combination --- the pair going to the concert is the same either way you list them.
$\binom{6}{2} = 15$.
\end{solution}
\end{problem}

\begin{problem}
In a school race with 10 runners, a gold, silver and bronze medallist will be awarded. How
many possible ways could those medals be awarded?
\begin{multipleChoice}
\choice[correct]{Permutation}
\choice{Combination}
\end{multipleChoice}

$\answer{720}$

\begin{solution}
Permutation --- the medals are different, so which of the three you get matters.
$P(10,3) = 10 \times 9 \times 8 = 720$.
\end{solution}
\end{problem}

\begin{problem}
A group of 25 people run a race, with the top 8 finishers advancing to the final. How many
possible groups advance?
\begin{multipleChoice}
\choice{Permutation}
\choice[correct]{Combination}
\end{multipleChoice}

$\answer{1081575}$

\begin{solution}
Combination --- all that matters is \emph{who} advances, not the order they finished in.
$\binom{25}{8} = 1{,}081{,}575$.

Contrast with the medals problem: there the finishing positions carried different
outcomes, so order mattered. Here the eight advancing runners are treated alike.
\end{solution}
\end{problem}

\begin{problem}
A group of 30 students must select a president, vice president, and treasurer. If each
position were filled randomly from among them, how many total leadership possibilities
are there?
\begin{multipleChoice}
\choice[correct]{Permutation}
\choice{Combination}
\end{multipleChoice}

$\answer{24360}$

\begin{solution}
Permutation --- the three roles are distinct, so it matters which person gets which.
$P(30,3) = 30 \times 29 \times 28 = 24{,}360$.
\end{solution}
\end{problem}

\begin{problem}
A corporation with 25 employees presents 3 employees each year with ``Employee of the
Year'' awards. In how many different ways can the employees be chosen?
\begin{multipleChoice}
\choice{Permutation}
\choice[correct]{Combination}
\end{multipleChoice}

$\answer{2300}$

\begin{solution}
Combination --- the three awards are identical, so only the set of winners matters.
$\binom{25}{3} = 2{,}300$.

Compare with the club officers: same idea of picking 3 people, but there the three roles
differed, making it a permutation.
\end{solution}
\end{problem}

\begin{problem}
The varsity track team is putting together a team to run a 3-stage relay race, where each
stage is different and run by a different team member. If there are 15 runners, how many
possibilities are there for making the relay team?
\begin{multipleChoice}
\choice[correct]{Permutation}
\choice{Combination}
\end{multipleChoice}

$\answer{2730}$

\begin{solution}
Permutation --- the stages are different, so assigning runner A to stage 1 and runner B to
stage 2 is a different relay team from the reverse. $P(15,3) = 15 \times 14 \times 13 =
2{,}730$.
\end{solution}
\end{problem}

\section*{Experimental vs Theoretical Probability}

An \textbf{experimental probability} uses actual recorded instances from ``playing out''
scenarios --- flipping a coin repeatedly and dividing the number of heads by the number
of flips.

A \textbf{theoretical probability} is determined from known information about the total
possible scenarios, without carrying any of them out --- saying there is a $1/2$ chance of
heads because a coin has two sides and one is heads.

\begin{problem}
Your chances of winning a lottery with 4,000,000 participants is
$\frac{1}{4{,}000{,}000} = 0.000025\%$.
\begin{multipleChoice}
\choice{Experimental}
\choice[correct]{Theoretical}
\end{multipleChoice}

\begin{solution}
Theoretical --- computed from the known structure of the lottery, with no lottery actually
run.
\end{solution}
\end{problem}

\begin{problem}
Since there are 4 aces in a standard deck of 52 cards, your chances of drawing an ace is
$\frac{4}{52} \approx 7.7\%$.
\begin{multipleChoice}
\choice{Experimental}
\choice[correct]{Theoretical}
\end{multipleChoice}

\begin{solution}
Theoretical --- derived from the known composition of the deck, with no cards drawn.
\end{solution}
\end{problem}

\begin{problem}
$11.1\%$ of US citizens were reported to be under the poverty line.
\begin{multipleChoice}
\choice[correct]{Experimental}
\choice{Theoretical}
\end{multipleChoice}

\begin{solution}
Experimental --- this comes from observing and counting actual people, not from any
theory about how many ought to be in poverty.
\end{solution}
\end{problem}

\begin{problem}
The likelihood of rolling an even number on a six-sided die is $\frac{3}{6} = 50\%$.
\begin{multipleChoice}
\choice{Experimental}
\choice[correct]{Theoretical}
\end{multipleChoice}

\begin{solution}
Theoretical --- three of the six faces are even, known without rolling anything.
\end{solution}
\end{problem}

\begin{problem}
In 2024, it snowed for 7 out of 28 days in February, so the likelihood of snow in February
was $\frac{7}{28} = 25\%$.
\begin{multipleChoice}
\choice[correct]{Experimental}
\choice{Theoretical}
\end{multipleChoice}

\begin{solution}
Experimental --- based on what was actually recorded that February. Note it also comes
from a single month, so it is a fairly thin basis for a general claim about February.
\end{solution}
\end{problem}

\begin{problem}
A farmer observes that 17 out of 20 seeds they planted sprouted, so the probability of a
seed sprouting is $\frac{17}{20} = 85\%$.
\begin{multipleChoice}
\choice[correct]{Experimental}
\choice{Theoretical}
\end{multipleChoice}

\begin{solution}
Experimental --- the farmer planted the seeds and counted the results. There is no theory
that predicts a sprouting rate in advance, which is exactly why it must be measured.
\end{solution}
\end{problem}

\section*{Independent vs Dependent Events}

Two events are \textbf{dependent} if one event's occurrence affects the other's. Pulling
an ace and then a two from a deck without replacement is dependent: removing the ace
changes what remains for the second draw.

Two events are \textbf{independent} if one event's occurrence does not affect the other. A
coin flip and a die roll are independent: how the coin lands does not change the die's
possible outcomes.

\begin{problem}
You roll two fair dice. The events are rolling a 3 on the first die, and rolling a sum
bigger than 7.
\begin{multipleChoice}
\choice[correct]{Dependent}
\choice{Independent}
\end{multipleChoice}

\begin{solution}
Dependent. The sum involves the first die, so knowing it landed on 3 changes the chance
the total exceeds 7 --- with a 3 showing you need better than a 4 on the second die,
whereas a 6 on the first would make it far easier.

This is the subtle case in the set: the two \emph{dice} are independent of each other, but
the two \emph{events} are not, because the second event is defined using the first die's
result.
\end{solution}
\end{problem}

\begin{problem}
You order a meal at a fast-food restaurant. The events are (A) you order a side of fries,
and (B) you order a drink.
\begin{multipleChoice}
\choice[correct]{Dependent}
\choice{Independent}
\end{multipleChoice}

\begin{solution}
Dependent. In practice these choices are related --- people ordering fries are more
likely to also order a drink, and combo meals bundle them together. Knowing someone
ordered fries does shift the probability that they ordered a drink.
\end{solution}
\end{problem}

\begin{problem}
A fair coin is tossed two times. The events are (A) first toss is a head and (B) second
toss is a tail.
\begin{multipleChoice}
\choice{Dependent}
\choice[correct]{Independent}
\end{multipleChoice}

\begin{solution}
Independent. The coin has no memory --- the second toss is $50/50$ regardless of the
first. Believing otherwise is the ``gambler's fallacy.''
\end{solution}
\end{problem}

\begin{problem}
You roll two fair dice. The events are rolling a 3 on the first die, and rolling an odd
number on the second die.
\begin{multipleChoice}
\choice{Dependent}
\choice[correct]{Independent}
\end{multipleChoice}

\begin{solution}
Independent. Each event concerns a different die, and neither is defined in terms of the
other. Compare this with the first problem in this section, where the second event
depended on \emph{both} dice.
\end{solution}
\end{problem}

\begin{problem}
Your friend Jean has blond hair. The events are (A) Jean has blond hair and (B) one of
Jean's parents has blond hair.
\begin{multipleChoice}
\choice[correct]{Dependent}
\choice{Independent}
\end{multipleChoice}

\begin{solution}
Dependent. Hair colour is inherited, so knowing Jean is blond raises the probability that
a parent is too.
\end{solution}
\end{problem}

\begin{problem}
You draw a card from a deck, then draw a second without replacing the first. The events
are (A) the first card is an Ace and (B) the second card is a seven.
\begin{multipleChoice}
\choice[correct]{Dependent}
\choice{Independent}
\end{multipleChoice}

\begin{solution}
Dependent. Without replacement, the first draw changes the deck: 51 cards remain for the
second draw instead of 52.

Note that if the first card \emph{were} replaced and the deck reshuffled, these events
would become independent. The dependence comes entirely from not replacing the card.
\end{solution}
\end{problem}

\end{document}
